% content/frontmatter-pages.tex
% -----------------------------------------------------------------------------
% SIGNATURE SHEETS, ABSTRACT, AND PREFACE
% -----------------------------------------------------------------------------
% Wording and layout of the three signature pages follow the faculty's 2026
% sample thesis book (Indonesian and English editions). Check them against the
% latest guide from your study program before printing.
%
% Which pages appear depends on \ThesisTarget:
%
%                              proposal   draft   semhas   full
%   LEMBAR PENGESAHAN              -        -        -       yes
%   LEMBAR PERSETUJUAN             -        -        -       yes
%   PERNYATAAN ORISINALITAS        -        -        -       yes
%   ABSTRAK                        -        -       yes      yes
%   KATA PENGANTAR                 -        -        -       yes
%
% A proposal is not defended and a supervision draft is not signed. The results
% seminar (semhas) needs the abstract but nothing that is signed.
%
% \ThesisBlankPage and \ThesisPageBreak (template.cls) insert a "sengaja
% dikosongkan" page in the screen edition and move to the next right-hand page
% in the print edition.
\ifThesisComplete

\ThesisBlankPage

\ifThesisFull
% ── LEMBAR PENGESAHAN: signed by both advisors, both examiners, and the dean.
% \BgThispage draws the faint UPN logo watermark on this page only.
\chapter*{\ApprovalSheetLabel}
\addcontentsline{toc}{chapter}{\ApprovalSheetLabel}
\BgThispage
\vspace{0.8cm}
\begin{center}{\bfseries\normalsize \ThesisTitleUpper\par}\end{center}
\vspace{0.2cm}
\begin{center}\ByLabel:\\ \StudentNameUpper\\ \NpmLabel. \StudentID\end{center}
\vspace{0.5cm}
\noindent \Lang{Telah dipertahankan di hadapan, dan disetujui oleh, Dewan Penguji Ujian \DocTypeName\ \StudyProgramPhrase, Fakultas Ilmu Komputer, Universitas Pembangunan Nasional Veteran Jawa Timur, pada tanggal \DefenseDate:}{Has been defended before, and accepted by, the Board of Assessors of the Thesis Examination of the \StudyProgramPhrase, Faculty of Computer Science, Universitas Pembangunan Nasional Veteran Jawa Timur, on \DefenseDate:}
\vspace{0.2cm}
\begin{center}\bfseries \ApprovedLabel\end{center}
\vspace{0.2cm}
\noindent\begin{tabular}{@{}p{0.67\linewidth} p{0.15\linewidth} p{0.25\linewidth}@{}}
\textbf{\underline{\AdvisorOne}} & \underline{\makebox[\linewidth]{\dotfill}} & \textbf{(\AdvisorOneRole)} \\
\AdvisorOneID & & \\[0.5cm]
\textbf{\underline{\AdvisorTwo}} & \underline{\makebox[\linewidth]{\dotfill}} & \textbf{(\AdvisorTwoRole)} \\
\AdvisorTwoID & & \\[0.5cm]
\textbf{\underline{\ExaminerOne}} & \underline{\makebox[\linewidth]{\dotfill}} & \textbf{(\ExaminerOneRole)} \\
\ExaminerOneID & & \\[0.5cm]
\textbf{\underline{\ExaminerTwo}} & \underline{\makebox[\linewidth]{\dotfill}} & \textbf{(\ExaminerTwoRole)} \\
\ExaminerTwoID & & \\
\end{tabular}
\vspace{0.8cm}
\begin{center}\begin{tabular}{@{}c@{}}\bfseries \EndorsedByLabel\\[1em]\bfseries \DeanRoleLabel\\[2cm]\bfseries {\underline{\DeanName}}\\\DeanID\end{tabular}\end{center}

\ThesisBlankPage

% ── LEMBAR PERSETUJUAN: the study-program coordinator approves the manuscript
% for the thesis examination.
\chapter*{\AcceptanceSheetLabel}
\addcontentsline{toc}{chapter}{\AcceptanceSheetLabel}
\BgThispage
\vspace{1cm}
\begin{center}{\bfseries\normalsize \ThesisTitleUpper\par}\end{center}
\vspace{0.2cm}
\begin{center}\ByLabel:\\ \StudentNameUpper\\ \NpmLabel. \StudentID\end{center}
\vspace{1cm}
\begin{center}\Lang{Disetujui untuk melanjutkan ke Ujian \DocTypeName}{Approved to proceed to the Thesis Examination}\end{center}
\vspace{6cm}
\begin{center}\begin{tabular}{@{}c@{}}\bfseries \ApprovedByLabel\\[1em]\bfseries \HeadProgramRoleLabel\\\bfseries \FacultyLabel\\[2cm]\bfseries {\underline{\HeadOfProgramName}}\\\HeadOfProgramID\end{tabular}\end{center}

\ThesisBlankPage

% ── PERNYATAAN ORISINALITAS
% Ask your study program whether a meterai (duty stamp) is required on the
% signature. The QR code links to \OriginalityQRLink; it is left out when that
% link is empty.
\chapter*{\PlagiarismStatementLabel}
\addcontentsline{toc}{chapter}{\PlagiarismStatementLabel}
\vspace{0.5cm}
\noindent \Lang{Saya yang bertanda tangan di bawah ini:}{I am the undersigned:}\par
\vspace{0.3cm}
\begin{tabular}{@{} l c l @{}}
\Lang{Nama}{Student Name} & : & \StudentName \\
\NpmLabel & : & \StudentID \\
\Lang{Gelar}{Degree Program} & : & \ProgramDegreeLabel \\
\Lang{Program Studi}{Study Program} & : & \StudyProgramName \\
\Lang{Fakultas}{Faculty} & : & \FacultyLabel \\
\end{tabular}
\vspace{0.5cm}

\noindent \Lang{Dengan ini saya menyatakan bahwa \DocTypeNameLower\ ini tidak mengandung bagian apa pun dari karya ilmiah lain yang telah diajukan untuk memperoleh gelar akademis di lembaga pendidikan tinggi mana pun. Selain itu, \DocTypeNameLower\ ini tidak mengandung karya atau pendapat yang sebelumnya ditulis atau diterbitkan oleh pihak lain, kecuali yang secara eksplisit dikutip dalam \DocTypeNameLower\ ini dan tercantum secara lengkap dalam daftar pustaka.}{Hereby declares that this \DocTypeNameLower\ contains no part of any other scientific work that has been submitted to obtain an academic degree at any higher education institution. Furthermore, it does not contain any work or opinions previously written or published by others, except for those which are explicitly cited in this thesis and listed completely in references.}

\vspace{\baselineskip}\noindent \Lang{Saya juga menyatakan bahwa dokumen ilmiah ini bebas dari unsur-unsur plagiarisme. Jika di masa mendatang ditemukan indikasi plagiarisme dalam \DocTypeNameLower\ ini, saya bersedia menerima sanksi sesuai dengan peraturan perundang-undangan yang berlaku.}{And I declare that this scientific document is free from elements of plagiarism. If in the future indications of plagiarism are found in this Thesis, I am willing to accept sanctions in accordance with the applicable laws and regulations.}

\vspace{\baselineskip}\noindent \Lang{Dengan demikian, saya membuat pernyataan ini tanpa paksaan dari siapa pun dan untuk digunakan sebagaimana mestinya.}{Thus, I made this statement without any coercion from anyone and to be used as it should.}
\par\vspace{1cm}
\noindent\begin{minipage}[b]{0.45\linewidth}\centering
  \ifthenelse{\equal{\OriginalityQRLink}{}}{}{\raisebox{\depth}{\qrcode[height=4.5cm]{\OriginalityQRLink}}}
\end{minipage}\hfill
\begin{minipage}[b]{0.5\linewidth}\raggedleft
  \Lang{Surabaya}{Surabaya}, \SignDate\\[0.3cm]\DeclarantLabel\\[2.5cm]\underline{\StudentNameUpper}\\\NpmLabel. \StudentID
\end{minipage}

\ThesisBlankPage
\fi

% ── ABSTRAK
% The identity table and the keywords line are single-spaced; the abstract text
% uses the body spacing (\ThesisBodySpacing), as in the sample book.
\chapter*{\Lang{ABSTRAK}{ABSTRACT}}
\addcontentsline{toc}{chapter}{\Lang{ABSTRAK}{ABSTRACT}}
\begingroup
\singlespacing
\noindent\begin{tabular}{@{} l c p{0.60\linewidth} @{}}
\StudentNameLabel\ / \NpmLabel & : & \StudentName\ / \StudentID \\
\ThesisTitleLabel & : & \ThesisTitle \\
\Lang{Dosen Pembimbing}{Advisor} & : & 1. \AdvisorOne \\
 & & 2. \AdvisorTwo \\
\end{tabular}
% \par closes the paragraph that holds the identity table first. Without it the
% \vspace runs in horizontal mode and the gap lands after the first line of the
% abstract instead of before it.
\par\vspace{0.5cm}
\ThesisBodySpacing
\input{content/\ThesisLanguage/abstract}
\par\vspace{0.5cm}
\singlespacing
\noindent {\bfseries \Lang{Kata Kunci}{Keywords}}: \Lang{\KeywordsID}{\KeywordsEN}
\par
\endgroup

\ifThesisFull
\ThesisBlankPage

% ── KATA PENGANTAR
\begingroup
\ThesisBodySpacing
\chapter*{\PrefaceLabel}
\addcontentsline{toc}{chapter}{\PrefaceLabel}
\input{content/\ThesisLanguage/preface}
% Signature space. \needspace keeps the whole block on one page.
\par\needspace{4.5cm}
\begin{flushright}
Surabaya, \SignDate\\[2.5cm]
\StudentName
\end{flushright}
\endgroup
\fi
\ThesisPageBreak

\fi
